\documentclass[12pt,a4paper]{article}

% --- Преамбула ---
\usepackage[utf8]{inputenc}
\usepackage[T1]{fontenc}
\usepackage[russian]{babel}
\usepackage{amsmath,amssymb}
\usepackage{hyperref}
\usepackage[margin=2.5cm]{geometry}
\usepackage{enumitem}
\usepackage{setspace}
\usepackage{booktabs}
\onehalfspacing

\title{\textbf{Второе дополнение к документу}\\
\large «Ортогональность как фундаментальное свойство частицы в Теории Мод»}
\subtitle{Полная классификация типов реальности на основе знаков энергии, массы и ортогональности}
\author{}
\date{Версия 1.0}

\begin{document}
\maketitle

\section{Введение}
В основном документе было введено понятие ортогональности $O$ и установлена её связь с массой $m$ и энергией $E$ через уравнение $E + m + O = \text{const}$. В первом дополнении было показано, что отрицательная ортогональность приводит к антигравитации и естественным образом объясняет природу тёмной энергии.

В данном дополнении делается следующий логический шаг: рассматриваются все возможные комбинации знаков трёх параметров --- $E$, $m$ и $O$. Показано, что таких комбинаций восемь, и каждая из них соответствует своему типу реальности --- вселенной с уникальными физическими законами. Наша Вселенная ($+ + +$) является лишь одной из восьми возможных.

Документ является прямым продолжением основного документа и первого дополнения. Он использует все введённые ранее определения и результаты. Документ самодостаточен --- все необходимые понятия определены в тексте.

\section{Теоретические основы}
Настоящий документ опирается на следующие понятия, определённые в основном документе и первом дополнении:

\begin{itemize}
    \item \textbf{Частица} $P_i$ --- фундаментальный узел сети. Свойства определяются исключительно её отношениями с другими частицами.
    \item \textbf{Мода} $M_{A \to B}$ --- отображение из дискретного причинного пути $\Gamma_{AB}$ в вещественные числа (поперечный спектр). Длина любой моды бесконечна.
    \item \textbf{Собственная мода} $M_{A \to A}$ --- замкнутый путь (петля), начинающийся и заканчивающийся на частице $A$. Содержит всю историю взаимодействий частицы (causal history).
    \item \textbf{Энергия $E$} --- плотность рисунка поперечного спектра на локальной площади моды. Может быть как положительной, так и отрицательной.
    \item \textbf{Масса $m$} --- отношение локальной площади моды к ортогональности. Может быть как положительной, так и отрицательной.
    \item \textbf{Ортогональность $O$} --- угол поворота собственной моды частицы относительно её внешних мод. Может быть как положительной, так и отрицательной.
    \item \textbf{Гравитационная масса} --- определяется содержимым causal history (всей бесконечной длины моды). Не меняется при повороте моды.
    \item \textbf{Инертная масса} --- определяется ортогональностью (углом поворота здесь и сейчас). Меняется при повороте моды.
\end{itemize}

\textbf{Основное уравнение} (для покоящейся частицы):
\begin{equation}
    E + m + O = \text{const}
    \label{eq:sum}
\end{equation}
где константа своя для каждого семейства частиц.

\section{Отрицательная энергия в Теории Мод}

\subsection{Физический смысл}
Энергия $E$ в Теории Мод --- это плотность рисунка поперечного спектра на локальной площади моды. Положительная энергия ($E > 0$) означает, что на данном участке моды преобладают «пики» (положительные отклонения от равновесия). Отрицательная энергия ($E < 0$) означает, что преобладают «провалы» (отрицательные отклонения).

\subsection{Поведение частиц с $E < 0$}
Частица с отрицательной энергией:
\begin{itemize}
    \item Стремится поглотить энергию из окружения, чтобы скомпенсировать свои «провалы».
    \item При взаимодействии с обычной частицей ($E > 0$) их энергии могут скомпенсироваться, что приведёт к аннигиляции с выделением или поглощением энергии, в зависимости от соотношения модулей.
    \item Может существовать стабильно только в среде, где окружение также имеет отрицательную энергию (в другой вселенной с другими знаками параметров).
\end{itemize}

\subsection{Отрицательная энергия и вакуум}
В стандартной физике вакуум --- низшее энергетическое состояние, и отрицательная энергия невозможна. В Теории Мод вакуум --- это скомпенсированное состояние сети мод, где сумма всех $F$-значений по замкнутому пути равна нулю. Локальное нарушение компенсации может давать как положительные, так и отрицательные значения $E$. Поэтому $E < 0$ не запрещено фундаментально --- оно лишь означает, что частица является «провалом» в сети, а не «пиком».

\section{Полная классификация: восемь типов реальности}
Каждый из трёх параметров может принимать положительное или отрицательное значение. Это даёт $2^3 = 8$ возможных комбинаций. Каждая комбинация соответствует своему типу материи (и, в более широком смысле, своей вселенной).

\begin{table}[h]
    \centering
    \caption{Восемь типов реальности}
    \begin{tabular}{cccccl}
        \toprule
        Тип & $E$ & $m$ & $O$ & Название & Краткое описание \\
        \midrule
        I   & $+$ & $+$ & $+$ & Наша Вселенная & Обычная материя \\
        II  & $+$ & $+$ & $-$ & Тёмная энергия & Антигравитация, не кучкуется \\
        III & $+$ & $-$ & $+$ & «Убегающая» материя & Гравитационное отталкивание \\
        IV  & $+$ & $-$ & $-$ & Двойная экзотика & Отталкивание и по гравитации, и по инерции \\
        V   & $-$ & $+$ & $+$ & Отрицательная энергия & Поглощает энергию из окружения \\
        VI  & $-$ & $+$ & $-$ & Комбинированная экзотика & Отрицательная энергия + антигравитация \\
        VII & $-$ & $-$ & $+$ & Комбинированная экзотика II & Отрицательная энергия + отталкивание \\
        VIII& $-$ & $-$ & $-$ & Полная инверсия & Зеркальная вселенная \\
        \bottomrule
    \end{tabular}
    \label{tab:types}
\end{table}

\subsection{Тип I: $E > 0, m > 0, O > 0$ --- Наша Вселенная}
\begin{itemize}
    \item Обычная материя, из которой состоим мы и всё наблюдаемое.
    \item Положительная энергия (излучает, передаёт энергию).
    \item Положительная гравитация (притягивается к другой массе).
    \item Положительная инерция (ускоряется в направлении приложенной силы).
\end{itemize}

\subsection{Тип II: $E > 0, m > 0, O < 0$ --- Тёмная энергия}
\begin{itemize}
    \item Описан в первом дополнении.
    \item Положительная энергия, положительная гравитация, отрицательная инерция.
    \item Не кучкуется, равномерно распределяется, создаёт отрицательное давление.
    \item Отождествляется с тёмной энергией.
\end{itemize}

\subsection{Тип III: $E > 0, m < 0, O > 0$ --- «Убегающая» материя}
\begin{itemize}
    \item Положительная энергия, отрицательная гравитация, положительная инерция.
    \item Отталкивается от обычной материи гравитационно, но на внешний толчок реагирует нормально.
    \item Должна быть вытеснена из галактик и скоплений.
    \item Может существовать в виде разреженного фона в пустотах (войдах).
    \item Кандидат на объяснение крупномасштабных пустот в структуре Вселенной.
\end{itemize}

\subsection{Тип IV: $E > 0, m < 0, O < 0$ --- Двойная экзотика}
\begin{itemize}
    \item Положительная энергия, отрицательная гравитация, отрицательная инерция.
    \item Отталкивается от всего: и гравитационно, и инерционно.
    \item Самая «неуловимая» форма материи.
    \item Может быть ответственна за самые быстрые и масштабные процессы расширения.
\end{itemize}

\subsection{Тип V: $E < 0, m > 0, O > 0$ --- Отрицательная энергия, обычная материя}
\begin{itemize}
    \item Отрицательная энергия, но нормальная гравитация и инерция.
    \item Частица поглощает энергию из окружения, стремясь скомпенсироваться.
    \item При контакте с нашей материей ($E > 0$) происходит аннигиляция с выделением энергии (если $|E_{\text{отр}}| < E_{\text{пол}}$) или с поглощением (если $|E_{\text{отр}}| > E_{\text{пол}}$).
    \item Может быть кандидатом на объяснение некоторых типов квантовых флуктуаций вакуума.
\end{itemize}

\subsection{Тип VI: $E < 0, m > 0, O < 0$ --- Комбинированная экзотика}
\begin{itemize}
    \item Отрицательная энергия, положительная гравитация, отрицательная инерция.
    \item Притягивается гравитационно, но реагирует на силу в противоположную сторону.
    \item Поглощает энергию из окружения.
    \item Чрезвычайно нестабильна в нашей Вселенной.
\end{itemize}

\subsection{Тип VII: $E < 0, m < 0, O > 0$ --- Комбинированная экзотика II}
\begin{itemize}
    \item Отрицательная энергия, отрицательная гравитация, положительная инерция.
    \item Отталкивается гравитационно, но реагирует на силу нормально.
    \item Поглощает энергию из окружения.
\end{itemize}

\subsection{Тип VIII: $E < 0, m < 0, O < 0$ --- Полная инверсия}
\begin{itemize}
    \item Все три параметра отрицательны.
    \item Это «зеркальная» вселенная, физически неотличимая от нашей, но состоящая из антиподов.
    \item Частицы из этой вселенной, попав в нашу, будут вести себя как полная противоположность нашей материи: отталкиваться гравитационно, отталкиваться инерционно и поглощать энергию.
    \item При столкновении частицы из вселенной VIII с частицей из вселенной I происходит полная аннигиляция с нулевым суммарным выделением энергии (если модули равны).
\end{itemize}

\section{Следствия классификации}

\subsection{Мультиверс как множество секторов}
Восемь комбинаций знаков задают восемь секторов Мультиверса. Каждый сектор --- это своя вселенная (или набор вселенных) с уникальным набором физических законов. Наша --- лишь одна из восьми.

Секторы не изолированы полностью. Частицы могут переходить из одного сектора в другой при изменении знака соответствующего параметра. Это «перепрограммирование» на уровне целых вселенных.

\subsection{Симметрии}
\begin{itemize}
    \item \textbf{Смена знака $E$} переводит частицу в сектор с инвертированной энергией.
    \item \textbf{Смена знака $m$} меняет гравитацию (притяжение $\leftrightarrow$ отталкивание).
    \item \textbf{Смена знака $O$} меняет инерцию (нормальная $\leftrightarrow$ обратная).
    \item \textbf{Одновременная смена всех трёх знаков} (I $\leftrightarrow$ VIII) даёт вселенную, физически неотличимую от нашей, но состоящую из антиподов. Это глобальная симметрия Мультиверса.
\end{itemize}

\subsection{Аннигиляция между секторами}
При контакте частиц из разных секторов возможны различные сценарии:
\begin{itemize}
    \item \textbf{I + VIII} (полные антиподы): полная аннигиляция, суммарная энергия равна нулю.
    \item \textbf{I + II} (наша материя + тёмная энергия): аннигиляции не происходит (разные знаки только у $O$). Они сосуществуют, но взаимодействуют слабо.
    \item \textbf{I + III} (наша материя + убегающая): гравитационное отталкивание, но нормальное инерционное взаимодействие.
\end{itemize}

\section{Предсказания}
\begin{enumerate}[label=\arabic*.]
    \item \textbf{Существование «убегающей» материи (Тип III).} Она должна быть вытеснена из галактик и образовывать фон в войдах. Её можно обнаружить по аномальному гравитационному отталкиванию на больших масштабах.
    \item \textbf{Аннигиляция при контакте с отрицательной энергией (Тип V).} Если в нашей Вселенной существуют частицы с $E < 0$, они должны аннигилировать с обычной материей. Это может проявляться как вспышки излучения в пустотах.
    \item \textbf{Переходы между секторами.} При экстремальных условиях (сверхвысокие энергии, сильные поля) возможен переход частицы из одного сектора в другой. Это может быть обнаружено как рождение «неправильных» частиц в коллайдерах.
    \item \textbf{Шумовые паттерны.} Каждый из восьми типов материи должен иметь свой уникальный шумовой паттерн. Сравнение шума от разных источников может выявить присутствие материи других типов.
\end{enumerate}

\section{Заключение}
Рассмотрение всех возможных комбинаций знаков энергии $E$, массы $m$ и ортогональности $O$ приводит к полной классификации из восьми типов реальности. Наша Вселенная ($+ + +$) --- лишь один из секторов Мультиверса. Остальные семь типов материи могут существовать в других вселенных или в виде примесей в нашей. Тёмная энергия естественно возникает как Тип II ($+ + -$). Другие типы могут объяснять крупномасштабные пустоты, аномальные вспышки и другие необъяснённые явления.

Документ является вторым дополнением к основному документу «Ортогональность как фундаментальное свойство частицы в Теории Мод» и использует все его определения и результаты. Документ самодостаточен --- все необходимые понятия определены в тексте.

\end{document}